% Presentasi: P-DINA + lapisan naratif Gemini untuk diagnosis learning gap.
% Build: latexmk -pdf slides.tex   (atau: pdflatex slides.tex dua kali)
\documentclass[aspectratio=169,10pt]{beamer}
\usepackage[T1]{fontenc}
\usepackage[utf8]{inputenc}
\usepackage{lmodern}
\usepackage[indonesian]{babel}
\usepackage{amsmath,amssymb}
\usepackage{booktabs,tabularx,array,colortbl}
\usepackage{tikz}
\usetikzlibrary{arrows.meta,positioning,shapes.geometric,calc,fit,backgrounds}
\usepackage{pgfplots}
\usepackage{listings}
\usepackage[most]{tcolorbox}
% Shared colours, TikZ styles and pgfplots defaults (used by slides.tex and naskah.tex).
\definecolor{navy}{HTML}{1F3A5F}
\definecolor{accent}{HTML}{D9822B}
\definecolor{tealc}{HTML}{2A9D8F}
\definecolor{gapred}{HTML}{C8553D}
\definecolor{okgreen}{HTML}{3A8D5C}
\definecolor{midgray}{HTML}{8A94A6}
\definecolor{lightbg}{HTML}{EEF2F7}
\definecolor{insuf}{HTML}{C9CED8}

\pgfplotsset{
  compat=1.18,
  every axis/.append style={
    font=\scriptsize,
    axis line style={midgray},
    tick style={midgray},
    label style={font=\scriptsize},
    tick label style={font=\scriptsize},
    legend style={font=\scriptsize, draw=none, fill=none},
    grid style={midgray!25},
    every node near coord/.append style={text=black!80},
  },
}
% Indonesian number format everywhere (0,5 and 1.000)
\pgfkeys{/pgf/number format/.cd, use comma, 1000 sep={.}}

% ragged-right columns (avoid stretched inter-word spacing in narrow cells)
\newcolumntype{P}[1]{>{\raggedright\arraybackslash}p{#1}}
\newcolumntype{L}{>{\raggedright\arraybackslash}X}

\tikzset{
  box/.style={draw=navy, rounded corners=3pt, fill=white, align=center, font=\scriptsize,
              inner sep=4pt, line width=0.6pt},
  boxnavy/.style={box, fill=navy, text=white},
  boxacc/.style={box, draw=accent, fill=accent, text=white},
  boxteal/.style={box, draw=tealc, fill=tealc, text=white},
  boxsoft/.style={box, draw=midgray, fill=lightbg},
  arr/.style={-{Stealth[length=2.2mm]}, line width=0.8pt, draw=navy!80},
  lbl/.style={font=\tiny, text=black!70, align=center},
}

% height used by every pgfplots figure; set before \input{figs/...}
\newlength\figH
\setlength\figH{5cm}


% ------------------------------------------------------------------ theme
\usetheme{default}
\setbeamercolor{structure}{fg=navy}
\setbeamercolor{normal text}{fg=black!85}
\setbeamercolor{frametitle}{fg=white,bg=navy}
\setbeamercolor{title}{fg=navy}
\setbeamercolor{alerted text}{fg=accent!90!black}
\setbeamerfont{frametitle}{size=\large,series=\bfseries}
\setbeamertemplate{navigation symbols}{}
\setbeamertemplate{itemize items}[circle]
\makeatletter
\setbeamertemplate{frametitle}{%
  \nointerlineskip
  \begin{beamercolorbox}[wd=\paperwidth,ht=3.2ex,dp=1.4ex,leftskip=1.2em]{frametitle}%
    \usebeamerfont{frametitle}\insertframetitle
    \ifx\insertframesubtitle\@empty\else{\usebeamerfont{framesubtitle}\normalsize\mdseries\ --\ \insertframesubtitle}\fi
  \end{beamercolorbox}\vspace{1ex}}
\makeatother
\setbeamertemplate{footline}{%
  \hbox to \paperwidth{\hspace{1.2em}{\tiny\color{midgray}P-DINA + Gemini LLM untuk diagnosis learning gap}%
  \hfill{\tiny\color{midgray}\insertframenumber/\inserttotalframenumber}\hspace{1.2em}}\vspace{1.2ex}}
\setlength{\leftmargini}{1.2em}

\newcommand{\src}[1]{\par\vfill{\tiny\color{midgray}Sumber: #1}}
\newtcolorbox{infobox}[1][]{colback=lightbg, colframe=navy, boxrule=0.4pt, arc=2pt, left=4pt, right=4pt,
  top=2pt, bottom=2pt, fontupper=\small, #1}
\newtcolorbox{warnbox}[1][]{colback=accent!8, colframe=accent, boxrule=0.4pt, arc=2pt, left=4pt, right=4pt,
  top=2pt, bottom=2pt, fontupper=\small, #1}
\lstset{basicstyle=\ttfamily\tiny, breaklines=true, columns=fullflexible, keepspaces=true,
  frame=none, backgroundcolor=\color{lightbg}, xleftmargin=3pt, xrightmargin=3pt}

\title{Diagnosis Learning Gap Matematika dengan P-DINA\\ dan Laporan Naratif LLM}
\subtitle{Pengembangan dengan data EEDI NeurIPS 2020 Task 3/4}
\author{Muhammad Fariq Faqih}
\date{\today}

\begin{document}

% ================================================================== title
\begin{frame}[plain]
  \vfill
  {\usebeamerfont{title}\usebeamercolor[fg]{title}\Large\bfseries\inserttitle\par}
  \vspace{0.6em}
  {\color{black!70}\insertsubtitle\par}
  \vspace{1.2em}
  {\color{accent}\rule{0.35\linewidth}{1.2pt}}\par
  \vspace{0.8em}
  {\small\insertauthor\par}
  {\small\color{midgray}\insertdate\par}
  \vfill
  {\tiny\color{midgray}Semua angka berasal dari artefak hasil run (\texttt{artifacts/}). EEDI adalah data pengembangan
  (mayoritas siswa SMP di UK), bukan validasi pada kelas SMA penelitian.}
\end{frame}

\begin{frame}{Agenda}
  \begin{columns}[T]
    \column{0.48\linewidth}
    \begin{enumerate}\small
      \item Posisi model di pipeline
      \item Pipeline end-to-end
      \item P-DINA: probabilistic DINA + variational EM
      \item Parameter awal dan proses training
      \item Data, evaluasi performa, dan alasan P-DINA
      \item Hyperparameter tuning dan kebijakan diagnosis
    \end{enumerate}
    \column{0.48\linewidth}
    \begin{enumerate}\small\setcounter{enumi}{6}
      \item Penyimpanan model (checkpoint)
      \item Pemilihan LLM dan spesifikasi
      \item Konfigurasi model
      \item System prompt dan pengaman
      \item Output sistem AI secara keseluruhan
      \item Kesimpulan, keterbatasan, langkah berikutnya
    \end{enumerate}
  \end{columns}
\end{frame}

% ================================================================== 1. position
\begin{frame}{Posisi tahap model di pipeline}
  \centering
  \resizebox{0.92\linewidth}{!}{% Where the two models sit: P-DINA computes, Gemini only communicates.
\begin{tikzpicture}[node distance=5mm and 6mm]
  \node[boxsoft, minimum width=2.1cm, minimum height=1.1cm] (resp) {Respons siswa\\[1pt]{\tiny benar / salah per soal}};
  \node[boxnavy, right=of resp, minimum width=2.4cm, minimum height=1.1cm] (pdina) {\textbf{P-DINA}\\[1pt]{\tiny menghitung P(mastery)}};
  \node[boxsoft, right=of pdina, minimum width=2.5cm, minimum height=1.1cm] (diag) {Diagnosis\\[1pt]{\tiny label, learning gap,}\\[-1pt]{\tiny rekomendasi R1--R5}};
  \node[boxsoft, right=of diag, minimum width=2.3cm, minimum height=1.1cm] (facts) {Fakta JSON\\[1pt]{\tiny agregat, ID pseudonim}};
  \node[boxacc, right=of facts, minimum width=2.4cm, minimum height=1.1cm] (llm) {\textbf{Gemini LLM}\\[1pt]{\tiny menyusun kalimat}};
  \node[boxsoft, below=8mm of llm, minimum width=2.4cm] (val) {Validator\\[1pt]{\tiny angka, heading, klaim}};
  \node[boxsoft, below=8mm of diag, minimum width=2.5cm] (app) {Aplikasi Streamlit\\[1pt]{\tiny tabel, grafik, laporan}};

  \draw[arr] (resp) -- (pdina);
  \draw[arr] (pdina) -- (diag);
  \draw[arr] (diag) -- (facts);
  \draw[arr] (facts) -- (llm);
  \draw[arr] (llm) -- (val);
  \draw[arr] (val) -- node[lbl, above]{lolos / template} (app);
  \draw[arr] (diag) -- node[lbl, right]{langsung} (app);

  \begin{scope}[on background layer]
    \node[fill=navy!7, rounded corners, fit=(resp)(pdina)(diag), inner sep=6pt,
          label={[font=\tiny\bfseries, text=navy]above:Tahap komputasi (deterministik, dapat diaudit)}] {};
    \node[fill=accent!10, rounded corners, fit=(facts)(llm)(val), inner sep=6pt,
          label={[font=\tiny\bfseries, text=accent!80!black]above:Tahap komunikasi (opsional)}] {};
  \end{scope}
\end{tikzpicture}
}
  \vspace{0.4em}

  \footnotesize
  \begin{tabularx}{\linewidth}{@{}l L L@{}}
    \toprule
    & \textbf{P-DINA} & \textbf{Gemini LLM} \\
    \midrule
    Peran & \emph{menghitung} P(mastery) per siswa $\times$ konsep & \emph{mengomunikasikan}: laporan berbahasa Indonesia \\
    Input & respons benar/salah + Q-matrix & fakta JSON hasil P-DINA (ID dipseudonimkan) \\
    Wajib? & ya, inti sistem & opsional; tanpa API key tampil template \\
    Dilarang & -- & menghitung/mengubah angka, menambah konsep, klaim efektivitas \\
    \bottomrule
  \end{tabularx}
\end{frame}

% ================================================================== 2. pipeline
\begin{frame}{Pipeline end-to-end}
  \centering
  \resizebox{\linewidth}{!}{% End-to-end pipeline: six stages.
\begin{tikzpicture}[
  stage/.style={box, minimum width=2.15cm, minimum height=2.5cm, text width=2.0cm, anchor=north},
  num/.style={circle, fill=#1, text=white, font=\scriptsize\bfseries, inner sep=1.5pt, minimum size=5mm},
]
  \foreach \xpos/\col/\ttl/\desc [count=\n] in {
      0/navy/{Data EEDI}/{1.382.727 jawaban\\4.918 siswa\\948 soal, 57 konsep},
      2.45/navy/{Audit \& Q-matrix}/{skema, leakage,\\soal\,$\to$\,konsep,\\11 konsep target},
      4.9/navy/{Training P-DINA}/{variational EM\\split per siswa\\80/10/10},
      7.35/navy/{Diagnosis \& rekomendasi}/{P(mastery), label,\\learning gap,\\aturan R1--R5},
      9.8/accent/{Laporan naratif LLM}/{Gemini 3.1\\Flash-Lite +\\validator},
      12.25/tealc/{Aplikasi Streamlit}/{asesmen, gap,\\intervensi,\\pre/post}} {
    \node[stage, draw=\col] (s\n) at (\xpos,0) {\textbf{\ttl}\\[3pt]{\tiny\desc}};
    \node[num=\col] at (s\n.north) {\n};
  }
  \foreach \a/\b in {1/2,2/3,3/4,4/5,5/6} {
    \draw[arr] ($(s\a.east)+(0,-1.25)$) -- ($(s\b.west)+(0,-1.25)$);
  }
\end{tikzpicture}
}
  \vspace{0.8em}

  \begin{columns}[T]
    \column{0.5\linewidth}
    \begin{infobox}
      \textbf{Satu arah data:} tahap 1--4 menghasilkan checkpoint dan diagnosis yang dapat diaudit;
      tahap 5 hanya membaca hasilnya, tidak pernah menulis balik.
    \end{infobox}
    \column{0.46\linewidth}
    \small
    \begin{itemize}
      \item Train $\to$ fit, validation $\to$ seleksi, test $\to$ evaluasi \textbf{sekali}.
      \item Aplikasi tidak melatih ulang; hanya memuat checkpoint.
    \end{itemize}
  \end{columns}
\end{frame}

% ================================================================== 3. P-DINA
\begin{frame}{P-DINA: probabilistic DINA + variational EM}{jenis masalah}
  \begin{columns}[T]
    \column{0.5\linewidth}
    \small
    \textbf{Masalah:} dari pola jawaban benar/salah, perkirakan \alert{probabilitas siswa menguasai tiap konsep}.
    \begin{itemize}
      \item \textbf{Variabel laten:} penguasaan $\alpha_{ik}\in\{0,1\}$ tidak pernah diamati; tidak ada label ``jawaban benar''.
      \item \textbf{Multi-atribut:} 57 konsep sekaligus, terhubung lewat Q-matrix (948 soal $\times$ 57 konsep, dari pakar).
      \item \textbf{Data jarang:} tiap siswa hanya menjawab sebagian soal; yang tidak dijawab $\neq$ salah.
    \end{itemize}
    \vspace{0.3em}
    \begin{infobox}
      \textbf{Intuisi DINA:} siswa yang menguasai semua konsep sebuah soal tetap bisa salah (\emph{slip} $s$);
      yang belum menguasai bisa benar karena menebak (\emph{guess} $g$).
    \end{infobox}
    \column{0.47\linewidth}
    \small
    \textbf{Model formal} (siswa $i$, soal $j$, konsep $k$)
    \begin{align*}
      \eta_{ij} &= \textstyle\prod_k \alpha_{ik}^{\,q_{jk}} \\
      P(X_{ij}{=}1\mid\boldsymbol\alpha_i) &= (1-s_j)^{\eta_{ij}}\, g_j^{\,1-\eta_{ij}} \\
      \alpha_{ik} &\sim \mathrm{Bern}\big(\sigma(a_k\theta_i - b_k)\big),\ \theta_i\sim\mathcal N(0,1)
    \end{align*}
    \vspace{-0.6em}
    \begin{itemize}\footnotesize
      \item Prior \emph{higher-order}: kemampuan umum $\theta_i$ membuat konsep dengan sedikit bukti ``meminjam'' informasi dari konsep lain.
      \item Mean-field $q(\boldsymbol\alpha_i)=\prod_k\mathrm{Bern}(\pi_{ik})$ menghindari enumerasi $2^{57}\approx1{,}4\times10^{17}$ profil.
    \end{itemize}
  \end{columns}
\end{frame}

\begin{frame}{P-DINA: probabilistic DINA + variational EM}{keunggulan}
  \begin{columns}[T]
    \column{0.56\linewidth}
    \footnotesize
    \begin{tabularx}{\linewidth}{@{}>{\bfseries}l L@{}}
      \toprule
      Keunggulan & Penjelasan \\
      \midrule
      Interpretabel & tiap soal punya $g$ dan $s$; tiap konsep punya P(mastery) yang langsung bermakna bagi guru \\
      Ringan & $\pm$2.000 parameter; satu fit $\pm$30 detik di CPU laptop \\
      Missing $\neq$ salah & likelihood hanya menjumlah respons teramati (mask) \\
      Tahu kapan tidak tahu & bukti sedikit $\Rightarrow$ \emph{insufficient\_evidence}, bukan gap \\
      Terkalibrasi & ECE test 0,008 (Q-logistic 0,036) \\
      Siswa baru & cukup E-step dengan parameter beku; 30 siswa $\pm$0,4 detik \\
      Update eksak & untuk 934 dari 948 soal (satu konsep), E-step = posterior Bayes tepat \\
      \bottomrule
    \end{tabularx}
    \column{0.4\linewidth}
    \footnotesize
    \textbf{Contoh Q-matrix} (biner, dari pakar)\\[3pt]
    \begin{tabular}{@{}lccc@{}}
      \toprule
      Soal & Writing Expr. & Linear Eq. & \dots \\
      \midrule
      Q-101 & 1 & 0 & \dots \\
      Q-102 & 0 & 1 & \dots \\
      Q-103 & 1 & 1 & \dots \\
      \bottomrule
    \end{tabular}\\[4pt]
    Q-103 menuntut \emph{kedua} konsep (aturan konjungtif ``AND''). Di EEDI hanya 14 soal yang multi-konsep.
    \vspace{0.6em}
    \begin{warnbox}
      Kualitas diagnosis bergantung pada ketepatan Q-matrix: audit 22 soal fokus, 17 cocok, 5 parsial.
    \end{warnbox}
  \end{columns}
\end{frame}

% ================================================================== 4. init + training
\begin{frame}{Parameter awal dan proses training}
  \centering
  \resizebox{0.9\linewidth}{!}{% Variational EM training loop.
\begin{tikzpicture}[node distance=4mm and 5mm]
  \node[boxsoft, text width=2.3cm] (init) {\textbf{Inisialisasi}\\{\tiny $g_j{=}s_j{=}0{,}2$, $a_k{=}1$, $b_k{=}0$, $\theta_i{=}0$}};
  \node[boxnavy, right=of init, text width=2.5cm] (e) {\textbf{E-step}\\{\tiny perbarui $\pi_{ik}$ (keyakinan mastery) dan $\theta_i$ (Newton)}};
  \node[boxnavy, right=of e, text width=2.5cm] (m) {\textbf{M-step}\\{\tiny $g_j, s_j$ tertutup (MAP Beta); $a_k, b_k$ (Newton)}};
  \node[boxsoft, right=of m, text width=2.1cm] (elbo) {\textbf{Hitung ELBO}};
  \node[diamond, draw=navy, aspect=2.2, inner sep=1pt, font=\tiny, align=center, right=of elbo] (conv) {$\Delta\text{ELBO}/|\text{ELBO}|$\\$\le 10^{-6}$?};
  \node[boxacc, below=7mm of conv, text width=1.9cm] (stop) {\textbf{Checkpoint}\\{\tiny $g, s, a, b$}};
  \draw[arr] (init) -- (e);
  \draw[arr] (e) -- (m);
  \draw[arr] (m) -- (elbo);
  \draw[arr] (elbo) -- (conv);
  \draw[arr] (conv) -- node[lbl, right]{ya (atau 150 iterasi)} (stop);
  \draw[arr] (conv.north) -- ++(0,0.55) -| node[lbl, pos=0.25, above]{tidak} (e.north);
\end{tikzpicture}
}
  \vspace{0.3em}

  \begin{columns}[T]
    \column{0.54\linewidth}
    \scriptsize
    \begin{tabularx}{\linewidth}{@{}l l L@{}}
      \toprule
      Parameter & Nilai awal & Keterangan \\
      \midrule
      $g_j$, $s_j$ (948 soal) & 0,2 & prior $\mathrm{Beta}(1{+}0{,}2\kappa;\,1{+}0{,}8\kappa)$ \\
      $a_k$, $b_k$ (57 konsep) & 1 dan 0 & prior higher-order \\
      $\theta_i$ (3.934 siswa) & 0 & kemampuan umum, MAP (Newton) \\
      $\pi_{ik}$ & E-step pertama & keyakinan mastery \\
      batas & $g,s\in[0{,}001;\,0{,}7]$ & dan $g+s\le0{,}98$ \\
      \bottomrule
    \end{tabularx}
    \column{0.42\linewidth}
    \footnotesize
    \begin{itemize}
      \item \textbf{E-step:} update tertutup untuk soal satu-konsep (2 perkalian matriks); 14 soal multi-konsep pakai titik tetap dengan \emph{damping} 0,5.
      \item \textbf{M-step:} $g_j, s_j$ rumus tertutup (MAP Beta).
      \item \textbf{Henti:} perubahan relatif ELBO $\le10^{-6}$ atau 150 iterasi.
      \item \textbf{Hasil terpilih:} konvergen dalam \alert{44 iterasi}, ELBO akhir $-630.546$, \alert{29,8 detik} CPU.
    \end{itemize}
  \end{columns}
\end{frame}

\begin{frame}{Hasil training: parameter yang dipelajari}
  \begin{columns}[c]
    \column{0.58\linewidth}
    \setlength\figH{5.6cm}
    % Learned guess/slip distribution for 948 items; artifacts/model/item_parameters.csv (bins of 0.05)
\begin{tikzpicture}
\begin{axis}[
  ybar, width=\linewidth, height=\figH, bar width=3.2pt,
  xmin=-0.02, xmax=0.72, ymin=0,
  xlabel={nilai parameter}, ylabel={jumlah soal},
  ymajorgrids, legend style={at={(0.5,1.03)}, anchor=south, legend columns=2}, legend cell align=left,
  xticklabel style={/pgf/number format/fixed, /pgf/number format/precision=1},
]
  \addplot+[fill=navy, draw=none] coordinates {
    (0.025,3) (0.075,9) (0.125,38) (0.175,67) (0.225,129) (0.275,165) (0.325,162)
    (0.375,106) (0.425,81) (0.475,74) (0.525,46) (0.575,24) (0.625,21) (0.675,23)};
  \addlegendentry{guess $g$ (rata-rata 0,338)}
  \addplot+[fill=accent, draw=none] coordinates {
    (0.025,101) (0.075,141) (0.125,125) (0.175,86) (0.225,74) (0.275,79) (0.325,71)
    (0.375,57) (0.425,57) (0.475,46) (0.525,36) (0.575,25) (0.625,19) (0.675,31)};
  \addlegendentry{slip $s$ (rata-rata 0,253)}
\end{axis}
\end{tikzpicture}

    \column{0.4\linewidth}
    \footnotesize
    \begin{tabular}{@{}lccc@{}}
      \toprule
      & rata-rata & median & rentang \\
      \midrule
      $g$ & 0,338 & 0,317 & 0,017--0,700 \\
      $s$ & 0,253 & 0,217 & 0,002--0,700 \\
      $1-s-g$ & 0,409 & 0,417 & 0,020--0,976 \\
      \bottomrule
    \end{tabular}
    \vspace{0.6em}
    \begin{itemize}
      \item $g\approx0{,}33$ wajar untuk pilihan ganda 4 opsi (tebak acak 0,25).
      \item Daya beda sedang $\Rightarrow$ satu konsep butuh $\pm$5 respons agar stabil.
      \item 900 dari 948 soal punya $\ge$30 respons train (parameter andal).
      \item Rata-rata $a_k=3{,}67$: penguasaan antar-konsep sangat terkait kemampuan umum $\theta$.
    \end{itemize}
  \end{columns}
  \src{\texttt{artifacts/model/item\_parameters.csv}, \texttt{pdina\_config.json}}
\end{frame}

% ================================================================== 5. data + evaluation
\begin{frame}{Data dan protokol evaluasi}
  \begin{columns}[T]
    \column{0.5\linewidth}
    \footnotesize
    \textbf{Dataset hasil preprocessing (dipakai penuh, tidak di-sampling)}\\[3pt]
    \begin{tabular}{@{}lrrl@{}}
      \toprule
      Split & Siswa & Interaksi & Dipakai untuk \\
      \midrule
      train & 3.934 & 1.112.450 & fit P-DINA \& baseline \\
      validation & 492 & 129.242 & seleksi HP \& kebijakan \\
      test & 492 & 141.035 & evaluasi akhir, sekali \\
      \midrule
      total & 4.918 & 1.382.727 & \\
      \bottomrule
    \end{tabular}\\[4pt]
    \begin{itemize}
      \item Split \textbf{per siswa} (80/10/10, seed 42): siswa tidak muncul di dua split.
      \item Model hanya membaca \texttt{UserId}, \texttt{QuestionId}, \texttt{IsCorrect}; kolom rawan leakage
            (pass rate, difficulty, prior accuracy, \dots) dikecualikan.
      \item Rekonstruksi lokal identik dengan file Drive (train/val: 0 baris berbeda).
    \end{itemize}
    \column{0.46\linewidth}
    \footnotesize
    \textbf{Protokol held-out}
    \begin{enumerate}
      \item Pada tiap siswa val/test, 20\% respons disembunyikan (min.\ 5 tetap terlihat).
      \item Parameter soal dibekukan dari train.
      \item Mastery ditebak dari konteks (E-step).
      \item Respons tersembunyi diprediksi $\Rightarrow$ AUC, log-loss, Brier, ECE.
      \item CI 95\%: bootstrap per siswa, 200 ulangan.
    \end{enumerate}
    \begin{warnbox}
      Tidak ada ground truth mastery di EEDI: prediksi respons yang baik \emph{tidak} membuktikan diagnosis benar,
      sehingga dilengkapi konsistensi split-half dan simulasi.
    \end{warnbox}
  \end{columns}
\end{frame}

\begin{frame}{Evaluasi performa: prediksi respons pada test}
  \centering
  \setlength\figH{3.6cm}
  % Test ROC-AUC with bootstrap 95% CI (left) and ECE (right); artifacts/evaluation/test_metrics.csv
% Rows (bottom -> top): 1 item-mean, 2 concept-score, 3 Q-logistic, 4 P-DINA independent, 5 P-DINA main
\begin{tikzpicture}
\pgfplotsset{
  rowbars/.style={xbar, bar width=7pt, ymin=0.4, ymax=5.6, ytick={1,...,5}, xmajorgrids, height=\figH,
                  scaled x ticks=false},
  ci/.style={draw=none, bar shift=0pt, fill=#1,
             error bars/.cd, x dir=both, x explicit, error bar style={black!70}},
}
\begin{axis}[
  name=auc, rowbars, scale only axis, width=0.42\linewidth,
  xmin=0.62, xmax=0.80,
  yticklabels={Item-mean, Concept-score, Q-logistic, P-DINA independen, \textbf{P-DINA (utama)}},
  xlabel={ROC-AUC test [CI 95\%]},
  xticklabel style={/pgf/number format/fixed, /pgf/number format/precision=2},
]
  \addplot+[ci=midgray]  coordinates {(0.6722,1) += (0.0095,0) -= (0.0114,0)};
  \addplot+[ci=midgray]  coordinates {(0.7182,2) += (0.0148,0) -= (0.0114,0)};
  \addplot+[ci=tealc]    coordinates {(0.7746,3) += (0.0116,0) -= (0.0089,0)};
  \addplot+[ci=navy!45]  coordinates {(0.7613,4) += (0.0102,0) -= (0.0099,0)};
  \addplot+[ci=navy]     coordinates {(0.7684,5) += (0.0110,0) -= (0.0097,0)};
\end{axis}
\begin{axis}[
  at={(auc.east)}, anchor=west, xshift=0.5cm, rowbars, scale only axis, width=0.27\linewidth,
  xmin=0, xmax=0.045, yticklabels={},
  xlabel={ECE (makin kecil makin baik)},
  xtick={0,0.01,...,0.04},
  xticklabel style={/pgf/number format/fixed, /pgf/number format/precision=2},
  nodes near coords, point meta=explicit symbolic,
]
  \addplot+[draw=none, bar shift=0pt, fill=midgray] coordinates {(0.0052,1) [$0{,}005$]};
  \addplot+[draw=none, bar shift=0pt, fill=midgray] coordinates {(0.0165,2) [$0{,}016$]};
  \addplot+[draw=none, bar shift=0pt, fill=tealc]   coordinates {(0.0363,3) [$0{,}036$]};
  \addplot+[draw=none, bar shift=0pt, fill=navy!45] coordinates {(0.0063,4) [$0{,}006$]};
  \addplot+[draw=none, bar shift=0pt, fill=navy]    coordinates {(0.0077,5) [$0{,}008$]};
\end{axis}
\end{tikzpicture}


  \footnotesize
  \begin{tabular}{@{}lcccccc@{}}
    \toprule
    Model & ROC-AUC [CI 95\%] & Log-loss & Brier & Accuracy & F1 & ECE \\
    \midrule
    \rowcolor{navy!8}\textbf{P-DINA (utama)} & 0,768 [0,759--0,779] & 0,568 & 0,194 & 0,699 & 0,704 & \textbf{0,008} \\
    P-DINA prior independen & 0,761 [0,751--0,772] & 0,575 & 0,197 & 0,693 & 0,701 & 0,006 \\
    Q-logistic & 0,775 [0,766--0,786] & 0,567 & 0,193 & 0,704 & 0,726 & 0,036 \\
    Concept-score & 0,718 [0,707--0,733] & 0,611 & 0,213 & 0,658 & 0,681 & 0,016 \\
    Item-mean & 0,672 [0,661--0,682] & 0,644 & 0,226 & 0,627 & 0,665 & 0,005 \\
    \bottomrule
  \end{tabular}
  \src{\texttt{artifacts/evaluation/test\_metrics.csv}; 492 siswa, 28.003 respons held-out}
\end{frame}

\begin{frame}{Evaluasi performa: kalibrasi dan kualitas diagnosis}
  \begin{columns}[T]
    \column{0.48\linewidth}
    \setlength\figH{5.6cm}
    % Reliability diagram on test held-out responses; artifacts/evaluation/calibration_test.csv
\begin{tikzpicture}
\begin{axis}[
  width=\linewidth, height=\figH, xmin=0, xmax=1, ymin=0, ymax=1,
  xlabel={rata-rata prediksi P(benar)}, ylabel={proporsi benar teramati},
  grid=major, legend pos=north west, legend cell align=left,
  xtick={0,0.2,...,1}, ytick={0,0.2,...,1},
  ticklabel style={/pgf/number format/fixed, /pgf/number format/precision=1},
]
  \addplot[dashed, midgray, domain=0:1] {x};
  \addlegendentry{kalibrasi sempurna}
  \addplot[navy, thick, mark=*, mark size=1.4pt] coordinates {
    (0.0754,0.0000) (0.1721,0.1808) (0.2588,0.2608) (0.3445,0.3536) (0.4470,0.4575)
    (0.5467,0.5436) (0.6497,0.6354) (0.7476,0.7628) (0.8541,0.8502) (0.9414,0.9465)};
  \addlegendentry{P-DINA (ECE 0,008)}
  \addplot[tealc, thick, mark=square*, mark size=1.3pt] coordinates {
    (0.0780,0.1940) (0.1572,0.2362) (0.2525,0.2615) (0.3517,0.3330) (0.4499,0.3950)
    (0.5500,0.5172) (0.6492,0.6279) (0.7500,0.7755) (0.8476,0.9073) (0.9346,0.9779)};
  \addlegendentry{Q-logistic (ECE 0,036)}
\end{axis}
\end{tikzpicture}

    \column{0.5\linewidth}
    \footnotesize
    \textbf{Kualitas diagnosis (proksi)}\\[3pt]
    \begin{tabularx}{\linewidth}{@{}L r@{}}
      \toprule
      Ukuran & Hasil \\
      \midrule
      Konsistensi split-half (test, 880 sel) & \textbf{0,876} \\
      Akurasi atribut, simulasi (ada / tanpa bukti) & 0,957 / 0,851 \\
      Recovery parameter, simulasi $r(g)$ / $r(s)$ & 0,965 / 0,982 \\
      \bottomrule
    \end{tabularx}\\[4pt]
    \textbf{Diagnosis 492 siswa test} (5.412 sel konsep target): 727 \emph{mastered}, 1.394 \emph{not\_mastered},
    3.291 \emph{insufficient\_evidence}; 322 siswa punya $\ge$1 gap.\\[4pt]
    \begin{infobox}
      Kurva P-DINA hampir berimpit dengan diagonal: P(benar) = 0,7 berarti $\pm$70\% benar.
      Q-logistic melenceng di ujung bawah dan atas.
    \end{infobox}
  \end{columns}
  \src{\texttt{calibration\_test.csv}, \texttt{run\_summary.json}}
\end{frame}

\begin{frame}{Alasan menggunakan P-DINA}
  \begin{columns}[T]
    \column{0.55\linewidth}
    \footnotesize
    \scriptsize\begin{tabular}{@{}l P{2.6cm} P{2.2cm}@{}}
      \toprule
      & \textbf{P-DINA} & \textbf{Q-logistic} \\
      \midrule
      AUC test & 0,768 & 0,775 \\
      Log-loss test & 0,568 & 0,567 \\
      Kalibrasi (ECE) & \textbf{0,008} & 0,036 \\
      Konsep ``menguasai'' & ya, P(mastery) per konsep & tidak (kompensatori) \\
      Parameter & guess \& slip per soal & bobot logistik \\
      \bottomrule
    \end{tabular}\\[3pt]
    {\scriptsize Selisih berpasangan (bootstrap per siswa): AUC $-0{,}007$ [$-0{,}011$; $-0{,}003$],
    log-loss $+0{,}001$ [$-0{,}003$; $+0{,}007$].}
    \begin{itemize}
      \item Akurasi prediksi praktis \textbf{setara}: log-loss tidak berbeda bermakna.
      \item Untuk keputusan pendidikan, probabilitas yang \textbf{terkalibrasi} dan \textbf{dapat dijelaskan}
            lebih penting daripada selisih AUC 0,007.
      \item Jauh lebih baik dari concept-score dan item-mean.
    \end{itemize}
    \column{0.42\linewidth}
    \footnotesize
    \textbf{Asesmen adaptif dari model yang sama}
    \setlength\figH{3.9cm}
    % Adaptive (mutual information) vs random item selection, replay on 300 test students;
% artifacts/evaluation/adaptive_replay_summary.csv
\begin{tikzpicture}
\begin{axis}[
  scale only axis, width=\dimexpr\linewidth-1.5cm\relax, height=\figH, xmin=4, xmax=21, ymin=0.75, ymax=0.93,
  xtick={5,10,15,20}, xlabel={panjang tes (jumlah soal)},
  ylabel={kesepakatan label vs.\ semua respons}, grid=major,
  legend pos=south east, legend cell align=left,
  yticklabel style={/pgf/number format/fixed, /pgf/number format/precision=2},
]
  \addplot[accent, thick, mark=*, mark size=1.5pt] coordinates {(5,0.789) (10,0.862) (15,0.892) (20,0.911)};
  \addlegendentry{adaptif (mutual information)}
  \addplot[midgray, thick, mark=square*, mark size=1.4pt] coordinates {(5,0.771) (10,0.824) (15,0.851) (20,0.871)};
  \addlegendentry{acak}
\end{axis}
\end{tikzpicture}

    Replay 300 siswa test: 10 soal adaptif menyamai label dari semua respons 86\% (acak 82\%).
  \end{columns}
\end{frame}

% ================================================================== 6. HP tuning
\begin{frame}{Hyperparameter tuning}
  \begin{columns}[T]
    \column{0.58\linewidth}
    \setlength\figH{6.2cm}
    % Validation held-out log-loss per configuration; artifacts/evaluation/hp_search_validation.csv
\begin{tikzpicture}
\begin{axis}[
  xbar, scale only axis, width=\dimexpr\linewidth-3.6cm\relax, height=\figH, bar width=5.5pt, scaled x ticks=false,
  xmin=0.570, xmax=0.5848, ymin=0.4, ymax=13.6, ytick={1,...,13},
  yticklabels={
    {indep., $\kappa{=}40$, cw}, {indep., $\kappa{=}2$, cw}, {indep., $\kappa{=}10$, cw},
    {indep., $\kappa{=}40$}, {indep., $\kappa{=}10$}, {indep., $\kappa{=}2$},
    {HO, $\kappa{=}40$, cw}, {HO, $\kappa{=}10$, cw}, {HO, $\kappa{=}2$, cw},
    {Q-logistic $\sigma{=}1$}, {HO, $\kappa{=}40$}, {HO, $\kappa{=}10$}, {\textbf{HO, $\kappa{=}2$ (terpilih)}}},
  xlabel={log-loss validation (respons held-out, makin kecil makin baik)}, xmajorgrids,
  xticklabel style={/pgf/number format/fixed, /pgf/number format/precision=3},
  nodes near coords, point meta=x,
  every node near coord/.append style={font=\tiny, /pgf/number format/fixed, /pgf/number format/precision=4},
]
  \addplot+[fill=midgray, draw=none, bar shift=0pt] coordinates
    {(0.5818,1) (0.5811,2) (0.5810,3) (0.5795,4) (0.5788,5) (0.5788,6)};
  \addplot+[fill=navy!45, draw=none, bar shift=0pt] coordinates
    {(0.5745,7) (0.5742,8) (0.5742,9) (0.5726,11) (0.5721,12)};
  \addplot+[fill=tealc, draw=none, bar shift=0pt] coordinates {(0.5735,10)};
  \addplot+[fill=accent, draw=none, bar shift=0pt] coordinates {(0.5721,13)};
\end{axis}
\end{tikzpicture}

    \column{0.4\linewidth}
    \footnotesize
    \textbf{Grid} (12 konfigurasi P-DINA + baseline):
    \begin{itemize}
      \item prior: \emph{independent} / \emph{higher-order} (HO)
      \item kekuatan prior Beta $\kappa\in\{2, 10, 40\}$
      \item class weighting (cw): ya / tidak
    \end{itemize}
    Kriteria: \textbf{log-loss} respons held-out validation.\\[4pt]
    \textbf{Temuan}
    \begin{itemize}
      \item HO konsisten lebih baik (0,572 vs 0,579).
      \item $\kappa$ hampir tidak berpengaruh: data per soal cukup.
      \item Class weighting memperburuk log-loss dan ECE di semua konfigurasi $\Rightarrow$ tidak dipakai.
    \end{itemize}
    \begin{infobox}
      Terpilih: \textbf{HO, $\kappa=2$, tanpa cw} $\Rightarrow$ prior $g,s\sim\mathrm{Beta}(1{,}4;\,2{,}6)$.
    \end{infobox}
  \end{columns}
  \src{\texttt{artifacts/evaluation/hp\_search\_validation.csv}}
\end{frame}

\begin{frame}{Hyperparameter tuning: kebijakan diagnosis}
  \begin{columns}[T]
    \column{0.48\linewidth}
    \setlength\figH{4.1cm}
    % Threshold tau selection on validation; artifacts/evaluation/threshold_selection_validation.csv
\begin{tikzpicture}
\begin{axis}[
  width=\linewidth, height=\figH, xmin=0.28, xmax=0.72, ymin=0.5850, ymax=0.5910,
  xlabel={threshold $\tau$}, ylabel={log-loss held-out (label keras)},
  grid=major, xtick={0.3,0.4,...,0.7},
  xticklabel style={/pgf/number format/fixed, /pgf/number format/precision=2},
  yticklabel style={/pgf/number format/fixed, /pgf/number format/precision=4},
  scaled y ticks=false,
]
  \fill[accent!12] (axis cs:0.425,0.5850) rectangle (axis cs:0.72,0.5910);
  \node[font=\tiny, text=accent!80!black, anchor=north] at (axis cs:0.575,0.5908) {dalam 1 SE dari yang terbaik};
  \addplot[navy, thick, mark=*, mark size=1.4pt] coordinates {
    (0.30,0.59024) (0.35,0.58909) (0.40,0.58826) (0.45,0.58661) (0.50,0.58607)
    (0.55,0.58612) (0.60,0.58607) (0.65,0.58583) (0.70,0.58609)};
  \addplot[only marks, mark=*, mark size=3pt, accent] coordinates {(0.50,0.58607)};
  \node[font=\tiny, text=accent!80!black, anchor=north east] at (axis cs:0.49,0.58600) {$\tau{=}0{,}5$ dipilih};
  \node[font=\tiny, anchor=north] at (axis cs:0.65,0.58575) {minimum mentah 0,65};
\end{axis}
\end{tikzpicture}

    \footnotesize
    \textbf{$\tau = 0{,}5$}: semua $\tau\in[0{,}45;0{,}70]$ dalam 1 SE dari minimum (0,65) $\Rightarrow$ aturan
    \emph{one-standard-error} memilih nilai netral 0,5.
    \column{0.48\linewidth}
    \setlength\figH{4.1cm}
    % Minimum evidence n_min from split-half classification consistency on validation (README_pdina.md, section 6)
\begin{tikzpicture}
\begin{axis}[
  ybar, width=\linewidth, height=\figH, bar width=9pt,
  ymin=0.78, ymax=0.90, xtick={1,2,3,4,5,6,8}, xmin=0.3, xmax=8.7,
  xlabel={$n_{\min}$ (respons per konsep, per paruh)}, ylabel={konsistensi split-half},
  ymajorgrids, nodes near coords, point meta=y,
  every node near coord/.append style={font=\tiny, /pgf/number format/fixed, /pgf/number format/precision=3},
  yticklabel style={/pgf/number format/fixed, /pgf/number format/precision=2},
]
  \addplot+[fill=navy!45, draw=none, bar shift=0pt] coordinates {(1,0.837) (2,0.825) (3,0.820) (4,0.839) (6,0.879) (8,0.881)};
  \addplot+[fill=accent, draw=none, bar shift=0pt] coordinates {(5,0.857)};
  \draw[gapred, dashed, thick] (axis cs:0.3,0.85) -- (axis cs:8.7,0.85);
  \node[font=\tiny, text=gapred, anchor=south] at (axis cs:7,0.85) {target 0,85};
\end{axis}
\end{tikzpicture}

    \footnotesize
    \textbf{$n_{\min}=5$}: nilai terkecil dengan konsistensi split-half $\ge0{,}85$. \textbf{$\delta=0$}: tanpa pita
    ketidakpastian target sudah tercapai.
  \end{columns}
  \vspace{0.3em}
  \footnotesize
  \centering
  \begin{tabular}{@{}llll@{}}
    \toprule
    \emph{insufficient\_evidence} & \emph{mastered} & \emph{not\_mastered} = learning gap & \emph{uncertain} \\
    respons $<n_{\min}$ & $\pi\ge\tau+\delta$ & $\pi<\tau-\delta$ & $|\pi-\tau|<\delta$ \\
    \bottomrule
  \end{tabular}
\end{frame}

% ================================================================== 7. storage
\begin{frame}{Penyimpanan model}
  \begin{columns}[T]
    \column{0.55\linewidth}
    \footnotesize
    \textbf{Checkpoint \texttt{artifacts/model/}} ($\pm$1,1 MB, tanpa GPU/framework DL)\\[3pt]
    \scriptsize
    \begin{tabular}{@{}>{\ttfamily}l P{3.3cm} r@{}}
      \toprule
      \normalfont File & Isi & Ukuran \\
      \midrule
      pdina\_params.npz & Q-matrix, $g$, $s$, $\lambda$, $a$, $b$, ID soal \& konsep & 459 KB \\
      pdina\_config.json & hyperparameter, 44 iterasi, ELBO & 1 KB \\
      diagnosis\_policy.json & $\tau$, $\delta$, $n_{\min}$, konsep target & 1 KB \\
      item\_parameters.csv & guess, slip, daya beda per soal & 202 KB \\
      item\_bank.csv, q\_matrix.csv & item bank \& Q-matrix & 331 KB \\
      concept\_metadata.csv, \dots & nama konsep, target, prasyarat & 23 KB \\
      intervention\_map\_resolved.csv & peta intervensi (placeholder) & 8 KB \\
      model\_card.json & ringkasan run, data, metrik & 17 KB \\
      \bottomrule
    \end{tabular}
    \column{0.42\linewidth}
    \footnotesize
    \textbf{Saat dipakai}
    \begin{enumerate}
      \item \texttt{CDMService.load()} memuat checkpoint dan memeriksa konsistensi ID soal, urutan konsep, konsep target.
      \item \texttt{diagnose()} menjalankan E-step saja (parameter beku).
      \item Aplikasi Streamlit \textbf{tidak} membaca dataset jutaan baris dan tidak melatih ulang.
    \end{enumerate}
    \textbf{Training ulang diperlukan bila:}
    \begin{itemize}
      \item ada soal baru (butuh $\ge$30 respons/soal),
      \item Q-matrix atau daftar konsep berubah,
      \item populasi berbeda (kelas SMA penelitian).
    \end{itemize}
  \end{columns}
\end{frame}

% ================================================================== 8. LLM choice
\begin{frame}{Pemilihan LLM dan spesifikasi}
  \begin{columns}[T]
    \column{0.55\linewidth}
    \footnotesize
    \textbf{Perbandingan yang sudah diukur}\\[3pt]
    \scriptsize
    \begin{tabularx}{\linewidth}{@{}l L L@{}}
      \toprule
      & SLM lokal (sebelumnya) & \textbf{Gemini 3.1 Flash-Lite} \\
      \midrule
      Model & Qwen3-4B-Instruct (Ollama) & \texttt{gemini-3.1-flash-lite} \\
      Tempat jalan & laptop, RTX 2050 4 GB & server Google (Gemini API) \\
      Waktu/laporan & guru $\pm$80--105 s & $\pm$6 s (30 siswa demo) \\
      Validator & lolos percobaan 1 & lolos percobaan 1 \\
      Beban laptop & VRAM penuh, baterai boros & hanya klien HTTP \\
      Kebutuhan & unduh 2,5 GB, server lokal & API key + internet \\
      \bottomrule
    \end{tabularx}
    \vspace{0.4em}
    \begin{itemize}
      \item Tugasnya hanya \emph{menulis ulang fakta} $\Rightarrow$ model kecil-cepat cukup; validator menjaga angka.
      \item Flash-Lite: varian paling ringan Gemini 3.1, cocok untuk banyak laporan pendek.
    \end{itemize}
    \column{0.42\linewidth}
    \footnotesize
    \textbf{Spesifikasi yang dipakai}\\[3pt]
    \begin{tabular}{@{}ll@{}}
      \toprule
      Akses & Gemini API (key AI Studio) \\
      SDK & \texttt{google-genai} 2.33.0 \\
      Model ID & \texttt{gemini-3.1-flash-lite} \\
      Mode & teks $\to$ teks, tanpa streaming \\
      Thinking & level \texttt{low} \\
      Request & 1 per laporan (+ repair) \\
      \bottomrule
    \end{tabular}\\[5pt]
    \textbf{Pemakaian terukur} (laporan guru, 30 siswa demo EEDI): 3.246 token input, 618 output, 119 thinking;
    5,97 detik; 1 percobaan.
    \begin{warnbox}
      Fakta (pseudonim, agregat) dikirim ke Google. Untuk data kelas nyata, gunakan tier berbayar dan cek ketentuan data.
    \end{warnbox}
  \end{columns}
\end{frame}

% ================================================================== 9. config
\begin{frame}{Konfigurasi model}
  \begin{columns}[T]
    \column{0.6\linewidth}
    \footnotesize
    \textbf{\texttt{configs/llm.json}}\\[3pt]
    \scriptsize
    \begin{tabular}{@{}>{\ttfamily}l >{\ttfamily}l P{3.3cm}@{}}
      \toprule
      \normalfont Kunci & \normalfont Nilai & Alasan \\
      \midrule
      model & gemini-3.1-flash-lite & versi stabil (preview juga tersedia) \\
      api\_key\_env & GEMINI\_API\_KEY & key tidak pernah ditulis di config \\
      env\_file & .env & cadangan bila env proses kosong \\
      temperature & 1.0 & rekomendasi Google untuk Gemini 3 \\
      top\_p & null & default API \\
      thinking\_level & low & latensi \& biaya rendah \\
      max\_tokens & 8192 & termasuk token thinking \\
      timeout\_sec & 120 & batas per request \\
      max\_retries & 2 & repair setelah gagal validasi \\
      prompt\_version & narrative-gemini-v1 & ikut kunci cache \\
      max\_class\_concepts & 11 & batas ukuran fakta (gap/rek.\ 5) \\
      max\_examples & 1 & contoh few-shot per request \\
      cache\_ttl\_hours & 168 & laporan sama $\Rightarrow$ $<$1 s \\
      \bottomrule
    \end{tabular}
    \column{0.37\linewidth}
    \footnotesize
    \textbf{Ketahanan}
    \begin{itemize}
      \item HTTP 429/5xx: retry dengan backoff 2 s / 8 s.
      \item 3 kegagalan beruntun: \emph{circuit breaker} 60 s.
      \item Output terpotong (\texttt{MAX\_TOKENS}) atau kosong ditolak.
      \item Tanpa key/internet: template langsung, aplikasi tidak crash.
    \end{itemize}
    \textbf{Audit log} per laporan: waktu, model, hash input, latensi, token, status validasi, fallback.
    \textbf{Tidak} menyimpan teks laporan, ID siswa, maupun key.
  \end{columns}
\end{frame}

% ================================================================== 10. system prompt
\begin{frame}[fragile]{System prompt}
  \begin{columns}[T]
    \column{0.64\linewidth}
    \footnotesize\textbf{Teks persis (audiens guru)} -- \texttt{cdacdm/narrative.py}
    \lstinputlisting[basicstyle=\ttfamily\fontsize{4.4pt}{5.1pt}\selectfont]{data/system_prompt_teacher.txt}
    \column{0.33\linewidth}
    \footnotesize
    \textbf{Struktur pesan ke Gemini}
    \begin{enumerate}
      \item \texttt{system\_instruction}: aturan di samping.
      \item 1 contoh few-shot (user $\to$ model) untuk audiens yang sama.
      \item \texttt{user}: ``FAKTA DIAGNOSIS (JSON)'' + instruksi menulis.
    \end{enumerate}
    \textbf{Audiens siswa:} sapa ``kamu'', tanpa data kelas, heading Ringkasan / Yang sudah dikuasai / Yang perlu
    dilatih / Langkah belajar / Catatan.
    \vspace{0.3em}
    \begin{infobox}
      Prompt hanya \emph{meminta}; kepatuhan \emph{diperiksa} oleh validator (slide berikut).
    \end{infobox}
  \end{columns}
\end{frame}

\begin{frame}{Pengaman di sekitar LLM}
  \centering
  \resizebox{0.95\linewidth}{!}{% LLM guardrail: validator, repair, deterministic fallback.
\begin{tikzpicture}[node distance=5mm and 8mm]
  \node[boxsoft, text width=1.9cm] (f) {Fakta JSON\\{\tiny dari P-DINA}};
  \node[boxsoft, right=of f, text width=2.2cm] (p) {Prompt\\{\tiny aturan + heading + 1 contoh}};
  \node[boxacc, right=of p, text width=1.9cm] (g) {\textbf{Gemini}\\{\tiny 1 request}};
  \node[boxnavy, right=of g, text width=1.9cm] (v) {\textbf{Validator}\\{\tiny hard fail}};
  \node[box, draw=okgreen, right=of v, text width=2.0cm] (ok) {\textbf{Tampilkan}\\{\tiny + cache disk}};
  \node[boxsoft, below=6mm of g, text width=2.4cm] (r) {Repair\\{\tiny catat pelanggaran, maks.\ 2$\times$}};
  \node[box, draw=gapred, below=6mm of ok, text width=2.0cm] (t) {\textbf{Template}\\{\tiny deterministik}};
  \draw[arr] (f) -- (p);
  \draw[arr] (p) -- (g);
  \draw[arr] (g) -- (v);
  \draw[arr] (v) -- node[lbl, above]{lolos} (ok);
  \draw[arr] (v) |- node[lbl, pos=0.25, right]{gagal} (r);
  \draw[arr] (r) -- (g);
  \draw[arr] (r.east) ++(0,-0.15) -- node[lbl, below, pos=0.6]{masih gagal / API mati} ($(t.west)+(0,-0.15)$);
\end{tikzpicture}
}
  \vspace{0.6em}

  \begin{columns}[T]
    \column{0.5\linewidth}
    \footnotesize
    \textbf{Validator (hard fail)}
    \begin{itemize}
      \item Setiap angka harus ada di fakta JSON, dan pada baris yang menyebut konsep harus milik konsep itu.
      \item Lima heading wajib lengkap dan berurutan; tanpa HTML/daftar bernomor.
      \item Rekomendasi placeholder wajib disebut ``perlu validasi guru''.
      \item Klaim terlarang: \emph{terbukti, dijamin, efektif meningkatkan, siswa malas}, \dots
      \item Konsep di luar fakta ditolak; panjang 120--900 kata.
    \end{itemize}
    \column{0.46\linewidth}
    \footnotesize
    \textbf{Privasi}
    \begin{itemize}
      \item Tidak ada matriks respons, nama, sekolah, atau \texttt{UserId} asli.
      \item Siswa disebut \texttt{Siswa-<sha256[:8]>}.
      \item Laporan guru hanya agregat kelas.
    \end{itemize}
    \textbf{Template deterministik} selalu lolos validator, jadi pengguna tidak pernah melihat teks yang belum
    divalidasi.
  \end{columns}
\end{frame}

% ================================================================== 11. output
\begin{frame}{Output sistem AI: diagnosis P-DINA}
  \begin{columns}[T]
    \column{0.5\linewidth}
    \footnotesize
    \textbf{Ringkasan kelas} (30 siswa demo EEDI, 8.653 respons)
    \setlength\figH{4.5cm}
    % Class summary of 30 EEDI demo students (facts sent to the LLM); slides/data/example_teacher_report.json
\begin{tikzpicture}
\begin{axis}[
  xbar stacked, scale only axis, width=\dimexpr\linewidth-4.0cm\relax, height=\figH, bar width=6pt,
  xmin=0, xmax=30, ytick={1,...,11}, enlarge y limits=0.06,
  yticklabels={Gradient Between Two Co-ordinates, Naming Co-ordinates in 2D, Factorising into a Single Bracket,
               Writing Expressions, Simplifying Expressions (Like Terms), Horizontal and Vertical Lines,
               Substitution into Formula, Solving Linear Inequalities, Midpoint Between Two Co-ordinates,
               Linear Sequences (nth term), Linear Equations},
  xlabel={jumlah siswa (dari 30)}, xmajorgrids,
  legend style={at={(0.5,1.03)}, anchor=south, legend columns=3, font=\tiny, /tikz/every even column/.append style={column sep=4pt}},
  y tick label style={font=\tiny},
]
  \addplot+[fill=gapred, draw=none] coordinates {(0,1) (8,2) (9,3) (12,4) (13,5) (5,6) (12,7) (4,8) (4,9) (6,10) (9,11)};
  \addlegendentry{learning gap (\emph{not\_mastered})}
  \addplot+[fill=okgreen, draw=none] coordinates {(0,1) (8,2) (7,3) (8,4) (8,5) (3,6) (7,7) (2,8) (2,9) (3,10) (4,11)};
  \addlegendentry{\emph{mastered}}
  \addplot+[fill=insuf, draw=none] coordinates {(30,1) (14,2) (14,3) (10,4) (9,5) (22,6) (11,7) (24,8) (24,9) (21,10) (17,11)};
  \addlegendentry{bukti kurang}
\end{axis}
\end{tikzpicture}

    \column{0.47\linewidth}
    \footnotesize
    \textbf{Profil satu siswa} (\texttt{Siswa-e7f6c011}, 7 dari 11 konsep)\\[3pt]
    \scriptsize
    \setlength\tabcolsep{4pt}
    \begin{tabular}{@{}l r r r l@{}}
      \toprule
      Konsep & bukti & benar & P(m.) & label \\
      \midrule
      Linear Sequences & 11 & 1 & 0,00 & gap \\
      Substitution into Formula & 22 & 8 & 0,00 & gap \\
      Factorising & 8 & 3 & 0,00 & gap \\
      Solving Linear Inequalities & 8 & 4 & 0,03 & gap \\
      Writing Expressions & 25 & 13 & 0,34 & gap \\
      Horiz.\ \& Vertical Lines & 4 & 0 & 0,00 & kurang \\
      Gradient & 1 & 1 & 0,47 & kurang \\
      \bottomrule
    \end{tabular}\\[4pt]
    \footnotesize
    ``kurang'' = bukti $<5$ respons: \textbf{tidak} dicap gap, walau probabilitasnya rendah.\\[4pt]
    \textbf{Rekomendasi R1--R5}: tindakan (remediasi / asesmen ulang / tambah bukti), level (seluruh kelas /
    kelompok / individual), aktivitas \& soal latihan dari item bank. Aktivitas masih \emph{placeholder}.
  \end{columns}
  \src{diagnosis \texttt{artifacts/demo/test\_split\_responses.csv} (30 siswa pertama), checkpoint EEDI}
\end{frame}

\begin{frame}{Output sistem AI: laporan naratif Gemini}
  \begin{columns}[T]
    \column{0.62\linewidth}
    \begin{tcolorbox}[colback=white, colframe=accent, boxrule=0.5pt, arc=2pt, left=5pt, right=5pt, top=3pt, bottom=3pt,
                      fontupper=\scriptsize, title={\scriptsize Laporan guru -- dihasilkan \texttt{gemini-3.1-flash-lite} (dipotong)},
                      coltitle=white, colbacktitle=accent]
      \textbf{Ringkasan kelas}\\
      Data mencakup 30 siswa yang diuji pada 11 konsep target dengan total 8653 respons. Estimasi model mendeteksi
      18 siswa memiliki learning gap berdasarkan ambang batas diagnostik yang ditetapkan.\\[3pt]
      \textbf{Temuan utama}
      \begin{itemize}
        \item \textbf{Simplifying Expressions by Collecting Like Terms} memiliki 13 siswa yang belum menguasai, dengan
              estimasi probabilitas penguasaan 35\%.
        \item \textbf{Substitution into Formula} menunjukkan 12 siswa belum menguasai, dengan estimasi probabilitas
              penguasaan 30\%.
        \item \textbf{Gradient Between Two Co-ordinates} mencatat 30 siswa dengan bukti tidak cukup, yang berarti
              memerlukan soal tambahan [\dots]
      \end{itemize}
      \textbf{Rekomendasi tindakan}
      \begin{itemize}
        \item \textbf{Simplifying Expressions by Collecting Like Terms} (seluruh kelas): Penjelasan ulang dengan contoh
              terbimbing, lalu latihan bertahap dengan umpan balik untuk 13 siswa. Rekomendasi ini adalah placeholder dan
              perlu validasi guru. [\dots]
      \end{itemize}
    \end{tcolorbox}
    \column{0.35\linewidth}
    \footnotesize
    \textbf{Metadata request}\\[3pt]
    \begin{tabular}{@{}ll@{}}
      \toprule
      Validator & lolos \\
      Percobaan & 1 \\
      Latensi & 5,97 s \\
      Token in/out/think & 3.246 / 618 / 119 \\
      Fallback & tidak \\
      \bottomrule
    \end{tabular}\\[5pt]
    Semua angka (30, 11, 8653, 18, 13, 35\%, \dots) berasal dari fakta P-DINA; LLM hanya merangkai kalimat.
    \begin{warnbox}
      Validator memeriksa angka, bukan makna pedagogis. Laporan adalah draf yang tetap ditinjau guru.
    \end{warnbox}
  \end{columns}
  \src{\texttt{slides/data/example\_teacher\_report.json} (dibuat 27-10-2026)}
\end{frame}

% ================================================================== 12. conclusion
\begin{frame}{Kesimpulan}
  \begin{enumerate}
    \item \textbf{P-DINA} (prior higher-order, variational EM) dilatih pada 1,1 juta respons train EEDI dalam
          $\pm$30 detik CPU dan menghasilkan P(mastery) per siswa $\times$ konsep yang \textbf{terkalibrasi} (ECE 0,008).
    \item Akurasi prediksi respons \textbf{setara} baseline terkuat (AUC 0,768 vs 0,775; log-loss tidak berbeda
          bermakna), dengan parameter yang dapat dijelaskan (guess, slip).
    \item Kebijakan diagnosis dipilih dari data validation: $\tau=0{,}5$, $n_{\min}=5$, $\delta=0$; konsistensi
          split-half test 0,876.
    \item \textbf{Gemini 3.1 Flash-Lite} hanya mengomunikasikan hasil: $\pm$6 detik per laporan, lolos validator pada
          percobaan pertama; tanpa API key sistem tetap berjalan dengan template.
    \item Seluruh sistem berjalan end-to-end di laptop (CPU) dan dikemas dalam aplikasi Streamlit.
  \end{enumerate}
\end{frame}

\begin{frame}{Keterbatasan dan langkah berikutnya}
  \begin{columns}[T]
    \column{0.5\linewidth}
    \footnotesize
    \textbf{Keterbatasan}
    \begin{itemize}
      \item EEDI = siswa SMP di UK; 11 konsep target hanya \emph{prasyarat} aljabar SMA, tanpa kalkulus.
      \item Tidak ada ground truth mastery; kualitas diagnosis hanya proksi.
      \item Q-matrix dianggap benar (5 dari 22 soal yang diaudit hanya cocok sebagian).
      \item Daya beda soal sedang: 3.291 dari 5.412 sel test masih \emph{insufficient\_evidence};
            \emph{Gradient} hanya 2 soal.
      \item Rekomendasi intervensi masih placeholder; tidak ada klaim efektivitas.
      \item LLM: baru beberapa laporan nyata; bergantung internet, kuota, dan pengiriman data ke Google.
    \end{itemize}
    \column{0.46\linewidth}
    \footnotesize
    \textbf{Langkah berikutnya}
    \begin{enumerate}
      \item Menyusun bank soal SMA (aljabar \& kalkulus) dengan Q-matrix yang ditinjau guru.
      \item Mengumpulkan respons kelas nyata dan melatih ulang (min.\ 30 respons per soal).
      \item Pre-test/post-test dengan kelompok pembanding.
      \item Guru memvalidasi peta intervensi.
      \item Mengevaluasi laporan LLM secara sistematis (\texttt{scripts/eval\_narrative.py}) dan menilainya bersama guru.
    \end{enumerate}
  \end{columns}
\end{frame}

\begin{frame}[plain]
  \centering\vfill
  {\Large\bfseries\color{navy} Terima kasih}\\[0.8em]
  {\small\color{midgray}Kode, artefak, dan naskah lengkap: \texttt{D:\textbackslash finalIL}, folder \texttt{slides/}}
  \vfill
\end{frame}

\end{document}
